\newpage
\section{Referencias}

\noindent Chapman, C., et al. (2022). Efficacy of Capture The Flag competitions in computer science education. \textit{Journal of Information Security Education}, 14(2), 11–25.

\vspace{0.3cm}

\noindent Cichonski, P., Millar, T., Grance, T., \& Scarfone, K. (2022). \textit{Computer Security Incident Handling Guide (NIST SP 800-61 Rev. 2)}. National Institute of Standards and Technology.

\vspace{0.3cm}

\noindent Deterding, S., Dixon, D., Khaled, R., \& Nacke, L. (2011). From game design elements to gamefulness: defining "gamification". En \textit{Proceedings of the 15th International Academic MindTrek Conference} (pp. 9–15). ACM.

\vspace{0.3cm}

\noindent Dörner, R., Göbel, S., Kickmeier-Rust, M., Masuch, M., \& Zweig, K. (2016). \textit{Entertainment Computing and Serious Games}. Springer.

\vspace{0.3cm}

\noindent González, L., Fernández, R., \& Castillo, M. (2021). Videojuego educativo para entrenamiento en ciberseguridad. \textit{Revista Iberoamericana de Seguridad Informática}, 8(2), 45–58.

\vspace{0.3cm}

\noindent Kapp, K. M. (2012). \textit{The Gamification of Learning and Instruction: Game-based Methods and Strategies for Training and Education}. Pfeiffer.

\vspace{0.3cm}

\noindent Liu, Y., Kim, H., \& Lee, D. (2020). Gamified cybersecurity training: A systematic review. \textit{Journal of Cyber Learning}, 5(4), 201–219.

\vspace{0.3cm}

\noindent Michael, D., \& Chen, S. (2005). \textit{Serious Games: Games That Educate, Train, and Inform}. Thomson Course Technology.

\vspace{0.3cm}

\noindent Mitnick, K., \& Simon, W. (2023). \textit{The Art of Deception: Controlling the Human Element of Security}. Wiley.

\vspace{0.3cm}

\noindent Pérez, V., \& Torres, J. (2023). \textit{Simuladores de red para entrenamiento en ciberseguridad: estudio en Bolivia} (Tesis de licenciatura). Universidad Mayor de San Andrés, La Paz.

\vspace{0.3cm}

\noindent Ramírez, P., Álvarez, J., \& López, S. (2022). Efectividad de la gamificación en la prevención de phishing. \textit{Congreso Latinoamericano de Ciberseguridad}, 1(1), 112–118.

\vspace{0.3cm}

\noindent Sánchez, J., García, F., \& Ruiz, M. (2021). Gamification design framework for cybersecurity training. \textit{IEEE Transactions on Education}, 64(3), 256–264.

\vspace{0.3cm}

\noindent Stallings, W. (2020). \textit{Cryptography and Network Security: Principles and Practice}. Pearson.

\vspace{0.3cm}

\noindent Vargas, E., Santos, D., \& Melo, F. (2019). Entrenamiento de ciberseguridad mediante realidad virtual. \textit{Journal of IT and Education}, 12(3), 77–86.

\vspace{0.3cm}

\noindent National Institute of Standards and Technology (NIST). (2020). \textit{NICE Framework (SP 800-181 Rev. 1): Workforce Framework for Cybersecurity}. Gaithersburg, MD: NIST.

\vspace{0.3cm}

\noindent Verizon Communications. (2025). \textit{2025 Data Breach Investigations Report}. New York, NY: Verizon.